\documentclass[11pt]{amsart}
\usepackage[localtheorems]{raeez-math-template}

\newtheorem{thm}{Theorem}[section]
\newtheorem{prop}[thm]{Proposition}
\newtheorem{lem}[thm]{Lemma}
\newtheorem{cor}[thm]{Corollary}
\theoremstyle{definition}
\newtheorem{defn}[thm]{Definition}
\newtheorem{rem}[thm]{Remark}

\newcommand{\kch}{\kappa_{\mathrm{ch}}}
\newcommand{\HH}{\mathrm{HH}}
\newcommand{\ChirHoch}{\mathrm{ChirHoch}}
\newcommand{\THH}{\mathrm{THH}}
\newcommand{\cC}{\mathcal C}
\newcommand{\cV}{\mathcal V}
\newcommand{\cZ}{\mathcal Z}
\newcommand{\Rep}{\mathrm{Rep}}
\newcommand{\Ran}{\mathrm{Ran}}
\newcommand{\RHom}{\mathrm{RHom}}
\newcommand{\End}{\mathrm{End}}
\newcommand{\C}{\mathbb C}
\newcommand{\Gm}{\mathbb G_m}
\newcommand{\Zderch}{Z^{\mathrm{der}}_{\mathrm{ch}}}
\newcommand{\SpCh}{\mathrm{SpCh}}
\newcommand{\PhiFA}{\Phi^{\mathrm{FA}}}
\newcommand{\Perf}{\mathrm{Perf}}
\newcommand{\CoHA}{\mathrm{CoHA}}
\newcommand{\Stab}{\mathrm{Stab}}
\newcommand{\Muk}{\Lambda_{\mathrm{Muk}}}
\newcommand{\KKconf}{K^{\kch}}

\title{Derived Chiral Centre and Koszul Complementarity}
\author{Raeez Lorgat}
\date{\today}

\begin{document}

\begin{abstract}
The boundary algebra produced by a Calabi-Yau threefold is an
$E_1$-chiral algebra.  Its braided object is not hidden inside that
boundary algebra.  It is the Drinfeld centre of the $E_1$ representation
category, and, under the Ben-Zvi-Francis-Nadler hypotheses, the category
of $E_2$-modules over the derived chiral centre.  The five objects
\[
  A,\qquad B(A),\qquad A^i,\qquad A^!,\qquad \Zderch(A)
\]
are distinct.  Bar-cobar inversion recovers $A$ from $B(A)$; Verdier
duality produces $A^!$ from $A^i$; Hochschild cochains produce the bulk.
Maulik-Okounkov stable envelopes give an $R$-matrix only on loci with
the required torus action and chamber data.  They do not, by themselves,
remove the compact centre-homotopy-colimit obstruction.
\end{abstract}

\maketitle

\tableofcontents

\section{Five Objects}
\label{sec:five-objects}

\begin{defn}[Boundary, bar, Koszul coalgebra, dual, and bulk]
\label{def:five-objects}
Let $A$ be an augmented $E_1$-chiral algebra in a fixed presentable stable
symmetric monoidal chiral ambient $\cV_{\mathrm{ch}}$.

\begin{enumerate}[label=\textup{(\alph*)}]
\item The \emph{boundary algebra} is $A$.
\item The \emph{bar coalgebra} is $B(A)$, an $E_1$-coassociative
coalgebra in the same completed chiral ambient.
\item The \emph{Koszul coalgebra} is
\[
  A^i := H^\bullet(B(A))
\]
when the cohomology coalgebra exists in the chosen completion.
\item The \emph{Verdier-Koszul dual algebra} is
\[
  A^! := \mathbb D_{\Ran}(A^i).
\]
In a finite strict model this is $H^\bullet(B(A))^\vee$ with the
transported multiplication.
\item The \emph{derived chiral centre} is
\[
  \Zderch(A)
  :=
  C^\bullet_{\mathrm{ch}}(A,A)
  \simeq
  \RHom_{A\otimes_{\mathrm{ch}} A^{\mathrm{op}}}(A,A)
\]
on the chiral Hochschild comparison surface.
\end{enumerate}
The bar-cobar counit
\[
  \Omega B(A) \longrightarrow A
\]
is an inversion statement.  It is not the definition of $A^!$.
\end{defn}

\begin{rem}[The centre is a bulk construction]
\label{rem:centre-is-bulk}
The Drinfeld centre and the derived chiral centre are centre, double, or
bulk constructions attached to $A$.  They do not change the native
operadic level of $A$.
\[
  \cZ\bigl(\Rep^{E_1}(A)\bigr)
  \simeq
  \Rep^{E_2}\bigl(\Zderch(A)\bigr)
\]
holds under the Ben-Zvi-Francis-Nadler Morita-dualisability and
Hochschild-comparison hypotheses.  The left side is a braided
representation category; the right side is its presentation by modules
over Hochschild cochains.  Neither side is the boundary algebra $A$.
\end{rem}

\section{Native Level at Dimension Three}
\label{sec:native-level}

\begin{prop}[The dimension-three output is native \(E_1\)]
\label{prop:d3-native-e1}
Let $\cC$ be a smooth proper Calabi-Yau category of dimension $3$ on the
witnessed object-level locus of the two-stage construction.  For a fixed
admissible specialisation datum $(\Sigma_2,C)$,
\[
  A_{\cC,C}
  :=
  \SpCh_{\Sigma_2,C}\bigl(\PhiFA_3(\cC)\bigr)
  \in E_1\text{-}\mathrm{ChirAlg}(C).
\]
The first naturally braided object attached to $A_{\cC,C}$ is
\[
  \cZ\bigl(\Rep^{E_1}(A_{\cC,C})\bigr),
\]
or, under the BZFN hypotheses,
\[
  \Rep^{E_2}\bigl(\Zderch(A_{\cC,C})\bigr).
\]
The $E_2$ structure belongs to this centre or module category, not to
$A_{\cC,C}$ itself.
\end{prop}

\begin{proof}
The two-stage factorisation fixes the native curve-level output by
\[
  n(d)=
  \begin{cases}
    \infty,& d=1,\\
    2,& d=2,\\
    1,& d\geq 3.
  \end{cases}
\]
At $d=3$, Stage 1 produces the holomorphic factorisation algebra
$\PhiFA_3(\cC)$ after the stated formality, framing, anomaly, completion,
and Hochschild-comparison witnesses.  Stage 2 integrates over
$\Sigma_2$ and restricts to the reference curve $C$, consuming the two
transverse complex directions and leaving an ordered $E_1$-chiral
algebra on $C$.  The Drinfeld centre is then the right adjoint to the
forgetful functor from braided half-braidings to ordered modules.  It is
a new construction on $\Rep^{E_1}(A_{\cC,C})$, not a native operation on
$A_{\cC,C}$.
\end{proof}

\begin{cor}[The \(R\)-matrix lives after the centre passage]
\label{cor:rmatrix-after-centre}
The nonsymmetric quantum-group braiding at $d=3$ is centre-level data.
For the local $\C^3$ chart,
\[
  \CoHA(\C^3)\cong Y^+(\widehat{\mathfrak{gl}}_1)
\]
is the positive half.  The full affine Yangian arises by Drinfeld
double or centre passage, and $\mathcal W_{1+\infty}$ appears through
the standard Fock or vacuum-module evaluation.  The positive half is
not the vertex algebra.
\end{cor}

\section{Homotopy Colimits and Stable Envelopes}
\label{sec:hocolim-stable-envelopes}

\begin{prop}[Centre and homotopy colimit do not commute in general]
\label{prop:centre-hocolim-gap}
For a diagram of $E_1$-chiral algebras $\{A_i\}_{i\in I}$, there is no
general equivalence
\[
  \cZ\!\left(\Rep^{E_1}\!\left(\operatorname*{hocolim}_{i\in I} A_i\right)\right)
  \simeq
  \operatorname*{hocolim}_{i\in I}
  \cZ\bigl(\Rep^{E_1}(A_i)\bigr).
\]
The obstruction is categorical: a half-braiding with the global
homotopy-colimit algebra is not the same datum as compatible
half-braidings on all local charts.
\end{prop}

\begin{proof}
A half-braiding for an object $M$ in the global centre is a coherent
system
\[
  M\otimes N \xrightarrow{\sim} N\otimes M
\]
natural in every module $N$ over the global algebra.  Local centres only
give such systems after restricting $N$ to each chart.  Descent from
local half-braidings to a global half-braiding requires higher coherence
on every overlap and higher overlap.  The $K3\times E$ non-locality
calculation records the witness
\[
  \rho_0=\frac{25}{26},\qquad
  \rho_1=\frac{1199}{1248}.
\]
These ratios measure the gap between the global centre and the homotopy
colimit of local centres in the computed degrees.  They witness
non-commutation; they do not say that the global centre fails to exist.
\end{proof}

\begin{prop}[Stable-envelope bypass: exact scope]
\label{prop:stable-envelope-scope}
The Maulik-Okounkov stable-envelope construction supplies an $R$-matrix
only after the following data are fixed:
\begin{enumerate}[label=\textup{(\roman*)}]
\item an algebraic torus $T$ acting on the relevant symplectic or
quasi-symplectic moduli space;
\item fixed loci and attracting chambers for this $T$-action;
\item equivariant cohomology or equivariant $K$-theory in which the
stable envelopes live;
\item compatibility with the Hall product, orientation data, completion,
and the intended representation category.
\end{enumerate}
Thus stable envelopes give a global route to a braiding on toric and
ADE/Kummer local models.  They do not give an unconditional global
bypass for a compact generic $K3\times E$.
\end{prop}

\begin{proof}
Stable envelopes are chamber-dependent equivariant correspondences.
Without a torus action there are no attracting chambers and no
stable-envelope operator
\[
  R=\Stab_C^{-1}\circ \Stab_{C'}.
\]
For a generic K3 surface $S$,
\[
  H^0(S,T_S)\cong H^0(S,\Omega^1_S)=H^{1,0}(S)=0,
\]
using the holomorphic symplectic form to identify $T_S$ with
$\Omega^1_S$.  Hence $\mathrm{Aut}(S)^\circ$ is trivial and no
two-dimensional algebraic torus acts on $S$.  The elliptic factor does
not supply the missing torus: $\mathrm{Aut}(E)^\circ\cong E$, an
abelian variety acting by translations, not $\Gm$.  The local
ADE/Kummer and toric cases are different because a toric chart supplies
the required $T=(\C^*)^2$ and chamber data.  The stable-envelope
statement is therefore restricted to those local models.
\end{proof}

\section{Scalar Complementarity Constants}
\label{sec:scalar-complementarity}

\begin{prop}[The constants that survive the local checks]
\label{prop:constants-survive}
The scalar complementarity constants used by the cross-volume
Theorem C lane are the following family-wise formulas:
\[
\begin{array}{c|c|c}
\text{family} & \text{duality datum} & \text{constant}\\
\hline
\text{free field/Kac-Moody} & \text{Koszul-rigid free pair} & 0\\
\text{Virasoro} & c\leftrightarrow 26-c & 13\\
\text{principal }W_3 & c\leftrightarrow 100-c,\quad \kch(W_3(c))=5c/6
  & (5/6)\cdot 100=250/3\\
\text{Bershadsky-Polyakov} & K_{\mathrm{BP}}=196,\quad \varrho_{\mathrm{BP}}=1/6
  & 196/6=98/3\\
\text{Mukai-enhanced K3} & \KKconf=2c_+(\Muk),\quad c_+(\Muk)=4
  & 8
\end{array}
\]
The constants $250/3$ and $98/3$ are respectively the principal $W_3$
and Bershadsky-Polyakov complementarity constants.  They do not come
from the Fake Monster lattice or from hyperbolic $E_{10}$.
\end{prop}

\begin{proof}
The $W_3$ line is immediate from the local Vol I formula
$\kch(W_3(c))=5c/6$ and the dual central charge relation
$c+c^!=100$:
\[
  \kch(W_3(c))+\kch(W_3(c^!))
  =
  \frac{5}{6}(c+c^!)
  =
  \frac{5}{6}\cdot 100
  =
  \frac{250}{3}.
\]
The Bershadsky-Polyakov line is the Vol I conductor formula
$\varrho_{\mathrm{BP}}K_{\mathrm{BP}}=(1/6)\cdot 196=98/3$.
The Mukai-enhanced K3 line uses the Vol III convention that
$\Muk=H^0(K3,\mathbb Z)\oplus H^2(K3,\mathbb Z)\oplus H^4(K3,\mathbb Z)$
has signature $(4,20)$, so $c_+(\Muk)=4$, and the conductor on that
face is $\KKconf=2c_+(\Muk)=8$.  The free and Virasoro values are the
standard Vol I complementarity values.  No additional Heegner
stabiliser factor is applied to the Mukai conductor unless a separate
normalisation theorem supplies it.
\end{proof}

\section{Consequences}
\label{sec:consequences}

\begin{thm}[Corrected derived-centre reading]
\label{thm:corrected-derived-centre-reading}
Let $A$ be a $d=3$ boundary chiral algebra produced on a witnessed
object-level locus of the CY-to-chiral construction.  Then:
\begin{enumerate}[label=\textup{(\roman*)}]
\item $A$ is native $E_1$.
\item $B(A)$ is the bar coalgebra, not a Swiss-cheese or centre object.
\item $\Omega B(A)\to A$ is bar-cobar inversion, not the definition of
$A^!$.
\item $A^i=H^\bullet(B(A))$ and $A^!=\mathbb D_{\Ran}(A^i)$ are the
Koszul coalgebra and Verdier-Koszul dual.
\item $\Zderch(A)=C^\bullet_{\mathrm{ch}}(A,A)$ is the derived chiral
centre.
\item Under BZFN hypotheses,
\[
  \cZ(\Rep^{E_1}(A))\simeq \Rep^{E_2}(\Zderch(A)).
\]
\item Maulik-Okounkov stable envelopes give an $R$-matrix only on the
equivariant loci where their hypotheses are present.
\item The compact global Hall-Drinfeld double remains a separate
construction, conditional on positive and negative halves, Cartan part,
pairing, completion, orientation data, and descent compatibility.
\end{enumerate}
\end{thm}

\begin{proof}
Items (i) through (vi) are the native-level and BZFN centre statements
of Sections~\ref{sec:five-objects} and~\ref{sec:native-level}.  Item
(vii) is Proposition~\ref{prop:stable-envelope-scope}.  Item (viii) is
the Hall-Drinfeld double scope in the Vol III quantum-group foundation:
a positive half is not yet a full double, and a centre presentation is
not a Hopf pairing with antipode and Cartan completion.
\end{proof}

\begin{rem}[Open proof obligations]
\label{rem:open-obligations}
Three proof obligations remain.
\begin{enumerate}[label=\textup{(\alph*)}]
\item A compact non-toric CY3 needs a completed hCS-to-Hall comparison
compatible with charge grading, orientation data, coproduct, and descent.
\item A global Hall-Drinfeld double needs the negative half, Cartan half,
non-degenerate pairing after quotienting the radical, completion, and
antipode.
\item A generic $K3\times E$ stable-envelope route would need a source
of equivariant chamber data absent from a generic K3 surface.
\end{enumerate}
\end{rem}

\end{document}
